\section{Collision coefficients and Gaussian identification}
\label{sec:gaussian}

The preceding section applies to every subsequential true-field law.
Fix one such limit $H$ and its uncentered moment functions $T_j$.
The plane covariance coefficient is
$v=(2\pi)^2c_*>0$, as in \cref{thm:plane-covariance}. We now show
that it determines all moments in $D$.

A related way to identify Gaussian pairings, by cancelling collision
singularities, appears in Kenyon's dimer-height proof
\cite[Lemma~3.1 and Proposition~3.2]{Kenyon2001}.  The singularities there
are computed from explicit determinant formulas.  Duminil-Copin,
Kozlowski, Lammers, and Manolescu use harmonicity and collision asymptotics
to identify higher full-plane increment correlations from the GFF
two-point function
\cite[Section~7.2]{DCKLM2026}.  Here separate harmonicity determines the
form of each singularity, while a stopped local coupling to the plane law
determines its coefficient.  The boundary traces proved in
\cref{sec:limits} then supply harmonic uniqueness in the bounded domain.

\begin{lemma}[The only isolated singularity is logarithmic]
\label{lem:log-singularity}
Fix a collision-free list $(y,\mathbf z)$ of interior points. In a
neighborhood of $x=y$ containing none of the points of $\mathbf z$,
\begin{equation}
 T_j(x,y,\mathbf z)
     =L(y,\mathbf z)\log\frac1{|x-y|}+R(x;y,\mathbf z),
 \label{eq:log-expansion}
\end{equation}
where $R$ extends harmonically across $x=y$ in the first variable.
The function $L$ is smooth on the parameter configuration space.
On compact subsets of that space one can choose a uniform radius
on which $R$, and each fixed order of its first-variable derivatives
on a smaller disk, are uniformly bounded.
\end{lemma}

\begin{proof}
By \cref{prop:moment-harmonicity,lem:no-contact}, the function of $x$
is harmonic on a punctured disk and has at most polylogarithmic
growth at its puncture. The polar Fourier expansion of a harmonic
function there has constant angular part $a+b\log r$ and, in angular
mode $n\ge1$, radial factors $r^n$ and $r^{-n}$. A coefficient of
$r^{-n}$ would grow faster than every power of $\log(1/r)$, so it
vanishes. The remaining positive modes extend harmonically to the
center. This proves \eqref{eq:log-expansion}. This argument can
also be obtained by taking Fourier coefficients on each concentric
circle and solving their radial differential equations; the growth
bound kills the negative coefficients one at a time.

There is a useful parameter formula that supplies the needed
uniformity. Choose a fixed radius $a>0$, small enough for all
$(y,\mathbf z)$ in a parameter neighborhood. Outward differentiation
in the first variable gives
\begin{equation}
 L(y,\mathbf z)
   =-\frac1{2\pi}\int_{|w-y|=a}
                      \partial_{n_w}T_j(w,y,\mathbf z)\,ds_w.
 \label{eq:circle-log-coefficient}
\end{equation}
The integral of the derivative of the removable remainder is zero,
and that of $\log(1/|w-y|)$ is $-2\pi$. Parametrize the circle by
$w=y+a e^{i\theta}$. The integrand is smooth in all parameters on
this fixed collision-free compact set, so $L$ is smooth. Subtract
the logarithmic term from the boundary values on this circle and
use the Poisson formula on its disk. It gives the extension $R$
and bounds all its first-variable derivatives on every fixed
smaller concentric disk, uniformly over compact parameter sets.
This proves in particular the continuity of $L$ needed below,
without inferring it from a nonuniform singular limit.
\end{proof}

\begin{proposition}[Identification of the collision coefficient]
\label{prop:collision-coefficient}
For $j\ge2$ and distinct $(y,\mathbf z)$,
\begin{equation}
                  L(y,\mathbf z)=\frac{v}{2\pi}T_{j-2}(\mathbf z).
 \label{eq:collision-coefficient}
\end{equation}
The convention is $T_0=1$; no centering is made in this identity.
\end{proposition}

\begin{proof}
\emph{A stopped comparison with the plane.}
We first establish precisely the mixing statement used to examine
the collision. Fix an interior center $y_0$, and fixed small
spectator balls away from it. For two smooth zero-integral tests
$\alpha,\beta$ in disjoint bounded disks, put
\[
 \alpha_r(x)=r^{-2}\alpha((x-y_0)/r),\qquad
 \beta_r(x)=r^{-2}\beta((x-y_0)/r).
\]
Their supports lie in $B(y_0,Cr)$. If $F_\delta$ is a product of
height averages on the spectator balls, then, with a plane copy
denoted by $H_\delta^{\mathrm{pl}}$,
\begin{align}
 &\limsup_{\delta\downarrow0}
 \left|\EE[H_\delta(\alpha_r)H_\delta(\beta_r)F_\delta]
    -\EE[H_\delta^{\mathrm{pl}}(\alpha_r)
                 H_\delta^{\mathrm{pl}}(\beta_r)]\,\EE F_\delta\right|
       \longrightarrow0\qquad(r\downarrow0).
 \label{eq:local-product-mixing}
\end{align}
The mesh limit is taken first at every fixed $r$.

Here is the coupling argument, including the absolute exterior
heights. Choose a fixed disk $B(y_0,R_0)\Subset D$ disjoint from
the spectators. In the spin representation of the true problem,
$W$ is fixed on the one connected boundary path. Imposing constancy
there and then its plus label leaves the $B$ marginal unchanged.
The double-cut stack estimate used in \cref{prop:plane-mixing}
therefore applies in this pin-free interior disk. Search from the
outside inward for a double $B$ circuit between radii of order
$r^{1/2}$ and $R_0$. It fails with probability at most a positive
power of $r$. On success the chosen circuit surrounds a disk of
radius comparable to $r^{1/2}$, and the stopping exploration has
not examined its interior.

The absolute heights at all spectator sites can be read in this
exterior. Choose the anchor on the prescribed boundary, with its
known height, and join it to the spectator balls by paths avoiding
$B(y_0,R_0)$. Such paths exist after reducing $R_0$ if needed:
paths in $D$ can be detoured around this interior disk. A finite
family of positive-width tubes and the spectator balls supplies
the same paths in sufficiently fine lattice approximations. The
edge increment formula, summed from the fixed anchor, gives every
height occurring in $F_\delta$ using only exterior spins. Affine
interpolation uses only the adjacent exterior sites. Thus revealing
these observations adds no interior information.

Conditional on the stopping circuit and all exterior spins,
including $W$, its interior is the free-$W$ double-cut law. Indeed
both $B$ values across every cut edge agree, so there is no constraint
on $W$ across it; the stopping rule has imposed no further interior
condition. The common worst-boundary coupling in the proof of
\cref{prop:plane-mixing}, applied inside this circuit, now matches
the spins in $B(y_0,Cr)$ to those of a plane copy with failure
probability at most a positive power of $r/r^{1/2}$. Its bound is
uniform over the exposed exterior. The conditional coupling can
be chosen with the plane law as its second marginal for every
exterior value. That marginal is consequently independent of the
exterior observations. On failure of the first search, choose
the plane copy independently as well. Altogether the matching
failure probability is at most $C r^\gamma$ for some $\gamma>0$,
in the limiting mesh regime.

On the matching event the two zero-total height smearings agree
with their plane versions. Equality of the local increments
determines heights only up to one additive constant, and the zero
integrals remove exactly that constant. The interpolation weights
have the same zero sum. All comparison statements can first be
implemented in finite free spin volumes and then passed to the
plane law.

Finally this coupling controls unbounded products. The true local
smearings have $L^q$ norms at most
$C_q(1+|\log r|)^{1/2}$ by \cref{prop:smearing-moments}; the plane
zero-total smearings have uniformly bounded $L^q$ norms by
\cref{prop:plane-moments}. The fixed spectator product has bounded
moments of every finite order. Applying H\"older with exponent
four to the two insertions, the spectator product, and the failure
indicator bounds the expectation of the product discrepancy by
\begin{equation}
                    C r^{\gamma/4}(1+|\log r|)=o(1).
 \label{eq:mixing-error-log}
\end{equation}
For the plane term the same bound holds with no logarithmic loss.
This proves \eqref{eq:local-product-mixing}. In particular we have
not used an arbitrary conditional spin law after revealing absolute
heights: those heights were first made measurable in the exterior
of an actual stopping cut.

\emph{Extracting the logarithmic coefficient.}
Choose $\alpha,\beta$ so that
\begin{equation}
 J=\iint\alpha(u)\beta(w)\log\frac1{|u-w|}\,du\,dw\ne0.
 \label{eq:nonzero-log-pairing}
\end{equation}
Such tests exist even in disjoint disks. For example, take the
first-coordinate derivatives of radial unit bumps about two points
on a horizontal line, with sufficiently small disjoint supports.
Integration by parts reduces their pairing to the corresponding
mixed derivative of the log kernel averaged over the two bumps.
That derivative is nonzero at their centers and keeps its sign
when the supports are small. Each derivative test has integral
zero.

Use fixed disjoint radial unit mollifiers at the points of
$\mathbf z$, and let $F$ be the product of the resulting height
averages; when $j=2$, take $F=1$. Separate harmonicity gives
$\EE F=T_{j-2}(\mathbf z)$. By \cref{thm:plane-covariance}, the
mesh limit of the plane pairing in
\eqref{eq:local-product-mixing} equals $vJ/(2\pi)$. The log
constant introduced by the rescaling vanishes since both tests
have integral zero. Taking the mesh limit and then $r\downarrow0$
in \eqref{eq:local-product-mixing} therefore yields
\begin{equation}
 \lim_{r\downarrow0}\EE[H(\alpha_r)H(\beta_r)F]
                  =\frac{vJ}{2\pi}T_{j-2}(\mathbf z).
 \label{eq:mixing-coefficient-limit}
\end{equation}

We compute the same limit from \cref{lem:log-singularity}. The two
shrinking supports are disjoint and avoid the spectators. Averaging
the spectator variables radially thus evaluates the moment function
at $\mathbf z$ exactly. The expression on the left of
\eqref{eq:mixing-coefficient-limit} is
\begin{align*}
 \iint\alpha(u)\beta(w)\bigg[
 &L(y_0+rw,\mathbf z)
       \left(\log\frac1r+\log\frac1{|u-w|}\right)\\
 &\hspace{17mm}+R(y_0+ru;y_0+rw,\mathbf z)\bigg]\,du\,dw.
\end{align*}
The term containing $\log(1/r)$ is exactly zero: for each fixed
$w$ its coefficient is independent of $u$, and $\int\alpha=0$.
This cancellation occurs before estimating the variation of $L$;
no modulus of continuity multiplied by $\log(1/r)$ is required.
Continuity of $L$ and the disjoint supports give convergence of
the remaining logarithmic term to $L(y_0,\mathbf z)J$.
For the remainder, subtract
$R(y_0;y_0+rw,\mathbf z)$ inside the $u$ integral. The subtracted
quantity again integrates to zero, and the uniform first-variable
derivative bound in \cref{lem:log-singularity} makes the remaining
integral $O(r)$. The limit is therefore $L(y_0,\mathbf z)J$.
Comparison with \eqref{eq:mixing-coefficient-limit} and $J\ne0$
proves \eqref{eq:collision-coefficient}.
\end{proof}

\begin{theorem}[Gaussian moment recursion]
\label{thm:gaussian-moments}
For distinct points $(x,y_1,\ldots,y_{j-1})$,
\begin{equation}
 T_j(x,\mathbf y)
   =m(x)T_{j-1}(\mathbf y)
      +v\sum_{i=1}^{j-1}G_D(x,y_i)
                    T_{j-2}(\mathbf y\setminus y_i).
 \label{eq:gaussian-recursion}
\end{equation}
For $j=1$ the assertion is $T_1=m$. These functions, as locally
integrable distributions, are the moment distributions of
$m+\sqrt v\,\Phi_D$, where $\Phi_D$ is the zero-boundary GFF with
covariance $G_D$. The limiting field $H$ has this law.
\end{theorem}

\begin{proof}
Fix the distinct spectator points. Subtract the right side of
\eqref{eq:gaussian-recursion} from its left side, and call the
result $U(x)$. Away from the spectator points it is harmonic in
$x$. The Green function convention is
\[
 G_D(x,y)=\frac1{2\pi}\log\frac1{|x-y|}
                     +\text{a function harmonic near }y.
\]
By \cref{prop:collision-coefficient}, the logarithmic coefficient
at each $y_i$ cancels. Lemma~\ref{lem:log-singularity} then shows
that $U$ extends harmonically across every $y_i$. Thus it is
harmonic on all of $D$.

The boundary estimate in \cref{lem:no-contact} bounds the moment
function as $x$ approaches any boundary point, including a mark.
The subtracted mean is bounded, and the Green terms tend to zero
at the boundary. Together with removability at the finitely many
interior singularities, this makes $U$ bounded in $D$. On each
open boundary arc, \eqref{eq:moment-boundary-trace} and the trace
of $m$ give $U\to0$. Bounded harmonic uniqueness with the two
exceptional points, proved in the boundary argument above, gives
$U=0$. More explicitly, compare $\pm U$ with
$\norm{U}_\infty$ times a harmonic extension of a boundary function
equal to one in a small neighborhood of each mark and zero outside
a slightly larger neighborhood, and let those neighborhoods shrink.
This proves the recursion.

Starting from $T_0=1$, the recursion uniquely gives the usual partial
pairing formula
\begin{equation}
 T_j(x_1,\ldots,x_j)
   =\sum_{\mathcal P}
       \prod_{\{a,b\}\in\mathcal P}vG_D(x_a,x_b)
       \prod_{c\notin\bigcup\mathcal P}m(x_c),
 \label{eq:partial-pairing-formula}
\end{equation}
where $\mathcal P$ ranges over collections of disjoint unordered
pairs in $\{1,\ldots,j\}$, including the empty collection.
Indeed the distinguished index is either a singleton, producing
the mean term, or is paired to exactly one other index, producing
one term of the sum in \eqref{eq:gaussian-recursion}. This is
also the uncentered Gaussian moment formula. Each displayed
product is locally integrable: the logarithmic Green singularities
are locally integrable, and its pairs have disjoint sets of
variables. By \cref{lem:no-contact}, agreement off diagonals is
agreement of the full moment distributions. It follows that all
joint smeared moments of $H$ equal those of $m+\sqrt v\,\Phi_D$.

We include the determinacy step. For a real test $f$, let
$Z=H(f)-\int mf$ and put
$s^2=v\iint f(x)G_D(x,y)f(y)\,dx\,dy$. The moment formula gives
$\EE Z^{2n}=(2n-1)!!\,s^{2n}$ and all odd moments zero. If
$s=0$ then $Z=0$ almost surely. Otherwise Cauchy--Schwarz gives
\[
 \EE|Z|^n\le (\EE Z^{2n})^{1/2}
                  \le s^n(2n)^{n/2}.
\]
Consequently $\sum_n |t|^n\EE|Z|^n/n!<\infty$ for every real
$t$. The exponential and characteristic-function series may be
integrated term by term, and the latter is
$\exp(-s^2t^2/2)$. Thus $Z$ is Gaussian. Applying this argument
to every real linear combination of a finite list of tests and
using the Cram\'er--Wold criterion proves the required joint
Gaussian law. Finally a countable dense collection of tests
determines the Borel law in the local Sobolev space of
\cref{prop:tightness}. Hence the random distribution itself has
the asserted law.
\end{proof}

\begin{corollary}[Full convergence and the odd-height normalization]
\label{cor:field-limit}
For every $t>1$,
\begin{equation}
 H_\delta\ \Longrightarrow\ m+\sqrt v\,\Phi_D,
 \qquad
 H_\delta-\EE H_\delta\ \Longrightarrow\ \sqrt v\,\Phi_D
 \quad\text{in }H^{-t}_{\mathrm{loc}}(D).
 \label{eq:integer-field-limit}
\end{equation}
All finite joint test moments converge. In particular
\begin{equation}
 \frac{h_\delta-\EE h_\delta}{\sigma}
       \ \Longrightarrow\ \Phi_D,
 \qquad
            \sigma=2\sqrt v=4\pi\sqrt{c_*}>0.
 \label{eq:odd-field-limit}
\end{equation}
The constant has the absolutely convergent finite-volume expression
of \cref{prop:normalization-series}; it is independent of $D,a,b$.
The uncentered odd-height limit has mean $2m+1$, the harmonic
extension of its prescribed values $+1,-1$.
\end{corollary}

\begin{proof}
Every subsequence admits a further limit by \cref{prop:tightness},
and \cref{thm:gaussian-moments} identifies every such limit with
the same law. Tightness and uniqueness therefore give full
convergence, with convergence of joint moments supplied by
uniform integrability.

In particular $\EE H_\delta(f)\to\int mf$ for each test $f$.
There is also deterministic local Sobolev convergence of the
means, which justifies subtracting the actual discrete expectation
in the stated topology. With a cutoff and Fourier basis as in
\eqref{eq:sobolev-bound}, each Fourier coefficient of
$\chi(\EE H_\delta-m)$ tends to zero and is bounded uniformly
in its index and in $\delta$. For $t>1$, the summable weights
$(1+|n|^2)^{-t}$ and dominated convergence show that the
$H^{-t}$ norm tends to zero. Apply this to the local exhaustion
and subtract this deterministic limit in the first convergence
of \eqref{eq:integer-field-limit}. Finally affine interpolation
preserves $h_\delta=2H_\delta+1$ exactly. Centering removes its
constant term and multiplies fluctuations by two, giving
\eqref{eq:odd-field-limit} with the stated normalization.
\end{proof}
